\chapter{Release Description}

\section{Intended use}
Please review the Statement of Work for specific intended use information.

\section{Functions implemented in this release}
\input{generated/tables/implemented-functions.tex}
   
\section{Functions not implemented in this release}
Monitor {\color{red}Battery Voltage State of Charge} serial data signal.


and
{\color{red}Battery Charge Disable Wakeup.}


Note: Planning to support these in MY23 SUSD release.



\section{Folder structure}
\includegraphics{images/SUSD_tree.png}

\section{Release Package Contents}

\subsection{ETAS Developed Software}
\input{generated/tables/released_software_ETAS.tex}

\subsection{Third-party Developed Software}
N/A

\section{Configuration Management Items (Documents only)}
\input{generated/tables/released_documents.tex}
   
\section{Target Environment Description}

\subsection{Hardware Target System and Components}
\begin{tabular}{ |l|l| }
	\hline
	{\bfseries Manufacturer} & {\bfseries}\\ \hline
	Model Number & ISOLAR-EVE\\ \hline
	Description & Virtual development and testing environment\\ \hline
\end{tabular}

\subsection{Software Prerequisites/Dependencies}
Refer to \file{SUM_SUSD_Design_Document.pdf}.

\subsection{Hardware and Software Tools}
\input{generated/tables/tools.tex}

Note: RTE command used for RTE generation

\begin{rtecommandline}
-nts --samples=memmap --os-define-osenv=RTAOS40 --exclusive-area-optimization=enable --os-output-param=all --have-64bit-int-types=1 --strict-unconnected-rport-check=off -tsl=256
\end{rtecommandline}

\subsection{Restrictions}
Please review to the Statement of Work for specific restrictions.

\section{Compile/Build Tools and Notes}

\subsection{Compiler/Build Notes}

This release of SUM\_SUSD has been tested using the following compiler:

\input{generated/tables/compilers.tex}

\section{Integration Notes}

\subsection{Integration into a Project}
There is no need for standard installation procedure for SUSD.
The folder itself contains all the necessary artifacts to integrate into a project.
Dropping the complete folder structure into a BSW configuration project will enable the user to configure SUSD and generate all the necessary artifacts.

\subsection{Configuring SUSD}
Once the complete package has been added to a BSW project, an additional module SUSD should appear.
From here you can open and begin configuring the module.
Specific descriptions of each configurable parameter available are available in the “Description” view of ISOLAR-A with BCT add-on when a parameter is selected.

\subsection{Generating Artifacts}
Generating the C/H/SWCD files from your configuration is similar to other BSW modules delivered from ETAS.
The Code Generation perspective of the ISOLAR-A with BCT add-on, should present the user with a dialog to select specific module within the project to generate.
Either a Partial Build or Full Build can be used, depending on the needs of the user.
In case of any configuration errors present in SUSD, the console window will output a description of any errors found during generation.
Further, when the user is configuring the module, the Problems Log will show any active errors in the configuration of SUSD.
A third visual aid is the appearance of red X’s in the configuration view over items which are invalid.
Simply hovering over these X’s will provide the user a description of the problem.
These are static validations.

\subsection{Configuring the RTE}
After the user has configured and generated the artifacts for SUSD, the output SWCD file \linebreak(\file{SUM_SUSD_SWCD.arxml}) can be used for the RTE configuration.
This is a complete SWC description per the AUTOSAR standard of all the users configured ports and interfaces required for the project specific configuration.
This SWCD is a standard ASW component and not a module description.
Therefore, there are no special considerations to be taken into account when configuring the RTE.
It is recommended that if specific changes are required to the SWCD, the user not manually edit this file, rather update and regenerate the configuration with the required changes.

In order to fulfill the requirement to store the Vehicle Power Mode in memory, AUTOSAR typed Per Instance Memory is used for storing this data in non-volatile memory.
Therefore, it is assumed the user will configure the NVM stack to support this required information.
Further, it is expected that during shutdown, the appropriate BSW services will write that specific NVM block to persistent storage.

\section{Installation}
This software does not require a standard installation. 
Please see Integration Notes above for further information on integrating SUSD into your project.

\subsection{Installation Instructions}
N/A

\section{Licensing}
Please review to the Statement of Work for specific licensing information.